\section{The shape of the law}\label{sec:shape}

This section reads two features of the shape off Theorems~\ref{thm:bulk}
to~\ref{thm:preparation}. The first is the global mode, where the golden ratio
enters, and the second is the interior minimum between the bulk and the plus end.

\subsection{The global mode}\label{sec:mode}

Let $P^{\max}_N=\max_{1\le k<N}P_N(k)$ be the largest interior bin and
$H_N=\sup_xh_N(x)$ the height of the profile. By
Proposition~\ref{prop:first-run} the two endpoint bins are
$P_N(N)=Q_\plus(N)$ and $P_N(0)=Q_\minus(N)$.

\begin{introtheorem}[the global mode]\label{thm:mode}
Assume (RV) and (SPAN), and let the preparation satisfy the bulk criterion
$B_n\Prob(D=n)\to0$.
\begin{enumerate}[(a)]
\item $P^{\max}_N=(H_N+o(1))/B_N$, every interior maximizer lies in $pN+O(B_N)$, and
along any subsequence on which $h_N$ converges uniformly the scaled maximizers
$(k-pN)/B_N$ converge to the set of maximizers of the limit profile.
\item If $\liminf_NB_NQ_\sgn(N)/H_N>1$ for some sign $\sgn$, the endpoint bin of
sign $\sgn$ exceeds every interior bin for all large $N$, so the global mode
of $P_N$ is an endpoint. If $\limsup_NB_NQ_\sgn(N)/H_N<1$ for both signs, the
global mode is interior for all large $N$.
\item If $Q_\sgn\in\RV_{-\rho_\sgn}$, the endpoint of sign $\sgn$ beats every
interior bin for all large $N$ when $\rho_\sgn<1/\alpha_*$, and loses to the
largest interior bin when $\rho_\sgn>1/\alpha_*$.
\item (The golden ratio.) For the stationary preparation,
$\rho_\sgn=\alpha_\sgn-1$. With equal indices $\alpha_\plus=\alpha_\minus=\alpha$,
the global mode is an endpoint for all large $N$ when $\alpha<\varphi=(1+\sqrt5)/2$
and is interior for all large $N$ when $\alpha>\varphi$. For two equal laws
$T_\plus(n)=T_\minus(n)=n^{-\alpha}$ and the stationary preparation,
$B_NQ_\sgn(N)/H_N\sim\mathcal C(\alpha)N^{1-\alpha+1/\alpha}$ for each sign, where
\begin{equation}\label{eq:Calpha}
 \mathcal C(\alpha)=\frac{\pi\bigl(K_\alpha\lvert\cos\frac{\pi\alpha}2\rvert\bigr)^{1/\alpha}}
 {4(\alpha-1)\,\Gamma(1+\frac1\alpha)\,\zeta(\alpha)^{1+1/\alpha}} ,
\end{equation}
and at $\alpha=\varphi$ the ratio tends to $\mathcal C(\varphi)=0.7761317\ldots<1$, so the
global mode is interior. For a fresh start of sign $\sgn_0$,
$Q_{\sgn_0}=T_{\sgn_0}$, $Q_{\osgn_0}=0$ and $\rho_{\sgn_0}=\alpha_{\sgn_0}>1>1/\alpha_*$,
so the global mode is interior for all large $N$ and every pair of indices.
\end{enumerate}
\end{introtheorem}

\begin{proof}
(a) Put $h_*=\inf_NH_N$, which is positive by Lemma~\ref{lem:profile}(d). By Theorem~\ref{thm:preparation}(b) the bulk criterion
gives
\begin{equation}\label{eq:mode-interior}
 \varepsilon_N=\sup_{1\le k<N}\bigl|B_NP_N(k)-h_N(x_{N,k})\bigr|\to0 .
\end{equation}
Hence $B_NP^{\max}_N=\max_{1\le k<N}h_N(x_{N,k})+O(\varepsilon_N)$, and
$\max_kh_N(x_{N,k})\le H_N$. For the lower bound, Lemma~\ref{lem:profile}(b) gives $K$
with $\sup_{|x|>K}h_N(x)<h_*/2$ for all $N$, so $H_N$ is the supremum of $h_N$ over
$[-K,K]$. For large $N$ the points $x_{N,k}$ with $1\le k<N$ cover $[-K,K]$ with mesh
$1/B_N$, since $pN/B_N\to\infty$ and $(1-p)N/B_N\to\infty$. The common Lipschitz
constant of Lemma~\ref{lem:profile}(a) gives $\max_kh_N(x_{N,k})\ge H_N-C_2/B_N$. Hence
$B_NP^{\max}_N=H_N+o(1)$.

Let $k_N$ be an interior maximizer. Then $h_N(x_{N,k_N})\ge
B_NP^{\max}_N-\varepsilon_N=H_N-o(1)>h_*/2$ for large $N$, so $|x_{N,k_N}|\le K$, which
is the statement $k_N\in pN+O(B_N)$. Along a subsequence with $h_N\to h$
uniformly we have $H_N\to\sup h$ and $h(x_{N,k_N})\ge H_N-o(1)$. Every limit
point $x$ of the bounded sequence $x_{N,k_N}$ satisfies $h(x)=\sup h$ by
continuity. So the distance from $x_{N,k_N}$ to the set of maximizers of $h$
tends to zero.

(b) If $B_NQ_\sgn(N)\ge(1+\eta)H_N$ for all large $N$, then by (a)
$B_NQ_\sgn(N)-B_NP^{\max}_N\ge\eta H_N-o(1)\ge\eta h_*-o(1)>0$, so the endpoint bin
of sign $\sgn$ exceeds every interior bin. If $B_NQ_\sgn(N)\le(1-\eta)H_N$ for
both signs and all large $N$, then (a) gives $B_NP^{\max}_N-B_NQ_\sgn(N)\ge\eta h_*-o(1)>0$
for both signs.

(c) By Lemma~\ref{lem:rv}(v), $B_N\in\RV_{1/\alpha_*}$. Hence
$B_NQ_\sgn(N)\in\RV_{1/\alpha_*-\rho_\sgn}$. A regularly varying function with a
positive index tends to infinity and one with a negative index tends to
zero~\cite[Proposition~1.5.1]{bgt1989}. Since $H_N$ is bounded above and below by
Lemma~\ref{lem:profile}(a) and~(d), part (b) applies.

(d) By Corollary~\ref{cor:stationary}(a) the stationary preparation satisfies the
bulk criterion and $Q_\sgn(N)\sim NT_\sgn(N)/(s(\alpha_\sgn-1))$, so
$\rho_\sgn=\alpha_\sgn-1$. With equal indices $\alpha_*=\alpha$, and $\alpha-1<1/\alpha$
if and only if $\alpha^2-\alpha-1<0$, that is, $\alpha<\varphi$. Part (c) then shows
that both endpoints beat the interior when $\alpha<\varphi$ and that both lose when
$\alpha>\varphi$. For unequal indices part (c) shows that the endpoint of sign $\sgn$
beats the interior if and only if $\alpha_\sgn-1<1/\alpha_*$, away from equality.

Now let $T_\plus(n)=T_\minus(n)=n^{-\alpha}$. Then $\mu_\sgn=\sum_{n\ge1}n^{-\alpha}=\zeta(\alpha)$,
$s=2\zeta(\alpha)$, $p=1/2$, $\theta_{\plus,N}=\theta_{\minus,N}=1$, and we may take
$a_\sgn(t)=t^{1/\alpha}$, so $B_N=(N/s)^{1/\alpha}$. The profile does not depend on $N$.
It is the density $h$ of $\frac12(Z_\plus-Z_\minus)$ with $Z_\plus$ and $Z_\minus$
independent copies of $Z_\alpha$. Its characteristic function is
\[
 \exp\bigl(K_\alpha(-it/2)^\alpha+K_\alpha(it/2)^\alpha\bigr)
 =\exp\bigl(-c_\alpha|t|^\alpha\bigr),\qquad c_\alpha=2^{1-\alpha}K_\alpha\bigl\lvert\cos\tfrac{\pi\alpha}2\bigr\rvert ,
\]
because $(-iy)^\alpha+(iy)^\alpha=2y^\alpha\cos\frac{\pi\alpha}2$ for $y>0$ and
$\cos\frac{\pi\alpha}2<0$. This function is positive and integrable, so by
Fourier inversion $h(x)\le h(0)$ for every $x$, and
\[
 H_N=h(0)=\frac1\pi\int_0^\infty e^{-c_\alpha t^\alpha}\,dt
 =\frac{\Gamma(1+\frac1\alpha)}{\pi\,c_\alpha^{1/\alpha}} .
\]
The stationary atom is $Q_\sgn(N)=s^{-1}\sum_{d\ge N}d^{-\alpha}\sim
N^{1-\alpha}/(2(\alpha-1)\zeta(\alpha))$. Multiplying,
\[
 \frac{B_NQ_\sgn(N)}{H_N}\sim\frac{N^{1/\alpha+1-\alpha}}{(2\zeta(\alpha))^{1/\alpha}\cdot2(\alpha-1)\zeta(\alpha)}
 \cdot\frac{\pi\,2^{(1-\alpha)/\alpha}\bigl(K_\alpha\lvert\cos\frac{\pi\alpha}2\rvert\bigr)^{1/\alpha}}{\Gamma(1+\frac1\alpha)}
 =\mathcal C(\alpha)\,N^{1-\alpha+1/\alpha},
\]
since $2^{(1-\alpha)/\alpha}2^{-1/\alpha}2^{-1}=2^{-2}$. At $\alpha=\varphi$ the
exponent vanishes, and $\alpha-1=1/\varphi$, $1+1/\alpha=\varphi$ and
$2-\alpha=\varphi^{-2}$, so $K_\varphi=\varphi\,\Gamma(\varphi^{-2})$ and
\[
 \mathcal C(\varphi)=\frac{\pi\varphi\bigl(\varphi\,\Gamma(\varphi^{-2})\lvert\cos\frac{\pi\varphi}2\rvert\bigr)^{1/\varphi}}
 {4\,\Gamma(\varphi)\,\zeta(\varphi)^{\varphi}}=0.7761317205\ldots .
\]
That $\mathcal C(\varphi)<1$ needs only four elementary bounds. Since $x^{-\varphi}$ is
convex, the trapezoid rule gives $\zeta(\varphi)\ge\frac12+1/(\varphi-1)=\varphi+\frac12$.
The minimum of $\Gamma$ on $(0,\infty)$ exceeds $0.8856$~\cite[Section~5.4(iii)]{dlmf},
so $\Gamma(\varphi)\ge0.8856$. Since $\Gamma\le1$ on $[1,2]$,
$\Gamma(\varphi^{-2})=\varphi^2\Gamma(1+\varphi^{-2})\le\varphi^2$. Finally
$\frac\pi2<\frac{\pi\varphi}2<\frac{5\pi}6$, so $\lvert\cos\frac{\pi\varphi}2\rvert\le\frac{\sqrt3}2$.
Together these give $\mathcal C(\varphi)\le\pi\varphi(\sqrt3\varphi^3/2)^{1/\varphi}/(4\cdot0.8856\,(\varphi+\frac12)^\varphi)=0.9513\ldots<1$,
and (b) shows that the global mode is interior.

For a fresh start of sign $\sgn_0$ we have $\Prob(D=n)=f_{\sgn_0}(n)\le T_{\sgn_0}(n)$,
and $B_nT_{\sgn_0}(n)\to0$ by Lemma~\ref{lem:rv}(v), so the bulk criterion holds.
Moreover $Q_{\sgn_0}=T_{\sgn_0}\in\RV_{-\alpha_{\sgn_0}}$ with
$\alpha_{\sgn_0}>1>1/\alpha_*$, and the other atom vanishes. Part (c) concludes.
\end{proof}

For two equal laws $H_N=h(0)$, so $\mathcal C(\alpha)$ also compares the endpoint
with the center bin. On the line $\rho_\sgn=1/\alpha_*$ the ratio
$B_NQ_\sgn(N)/H_N$ is slowly varying, the comparison depends on the slowly varying
factors, and we leave it open.

\subsection{The boundary layer at second order}\label{sec:second-order}

By Theorem~\ref{thm:bigjump} the law follows the renewal mass function $u_\minus$
in the boundary layer, and $u_\minus(j)\to1/\mu_\minus$ by~\eqref{eq:renewal-thm}.
If $u_\minus$ stays above its limit, the layer lies above the level
$T_\plus(N)/\mu_\minus$ at which it meets the big-jump zone, while the big-jump kernel
increases as $k$ decreases from $N$. To locate the interior minimum between the two
we need the boundary layer at second order, with an exact coefficient.

\begin{theorem}[second order in the boundary layer]\label{thm:second-order}
Assume (RV) and (LR) for the plus sign, and let the start be fresh with sign
$\plus$. Put
\begin{equation}\label{eq:Cj}
\begin{split}
 C(j)&=\alpha_\plus\Bigl[j\,u_\minus(j)+2\mu_\plus v_\minus(j)+\sum_{i=0}^{j-1}u_\minus(i)\Bigr]\\
 &=\alpha_\plus\Bigl[j\,u_\minus(j)+\mu_\plus v_\minus(j)
 +\sum_{i=0}^{j-1}u_\minus(i)\bigl(1+\mu_\plus u_\minus(j-i)\bigr)\Bigr].
\end{split}
\end{equation}
\begin{enumerate}[(a)]
\item For every fixed $j\ge0$,
\[
 \frac{P_N(N-j)}{T_\plus(N)}=u_\minus(j)+\frac{C(j)}{N}+o\Bigl(\frac1N\Bigr).
\]
\item For every $\eta\in(0,1)$,
\[
 \sup_{1\le j\le N^{1-\eta}}\frac Nj\,\Bigl|\frac{P_N(N-j)}{T_\plus(N)}-u_\minus(j)-\frac{C(j)}{N}\Bigr|\to0 .
\]
\end{enumerate}
If in addition (RV) and (SPAN) hold for the minus sign, then
$C(j)/j\to2\alpha_\plus s/\mu_\minus^2$ as $j\to\infty$.
\end{theorem}

In the second form of $C(j)$ the terms $j\,u_\minus(j)+\mu_\plus v_\minus(j)$ come
from a censored giant run that starts after $a$ slots, through
$T_\plus(N-a)/T_\plus(N)=1+\alpha_\plus a/N+o(a/N)$, and the sum comes from a complete
giant run followed by a fresh minus walk, through Lemma~\ref{lem:green}. We prove the
first form. The convergence of $C(j)/j$ is slow, of the order of $jT_\minus(j)$
by~\eqref{eq:second-renewal} below, and Section~\ref{sec:verify} reports the check.

The proof uses the final-run identity~\eqref{eq:final-explicit} and two renewal
identities besides~\eqref{eq:cover}. Write $u=u_\minus$ and $v=v_\minus$. For $j\ge1$,
\begin{align}
 \sum_{m=1}^{j}T_\minus(m)\,v(j-m)&=\sum_{i=1}^{j-1}u(i),\label{eq:cover-v}\\
 \sum_{i=0}^{j-1}u(i)\,u(j-i)&=v(j).\label{eq:conv-v}
\end{align}
For~\eqref{eq:cover-v}, $v(j-m)T_\minus(m)=\sum_rr\Prob(S^\minus_r=j-m,\,X^\minus_{r+1}\ge m)$,
so the sum is the expectation of the index $r$ of the run that covers level
$j$, minus one. That index minus one equals the number of $r\ge1$ with
$S^\minus_r\le j-1$, whose expectation is $\sum_{i=1}^{j-1}u(i)$. For
\eqref{eq:conv-v}, $\sum_{i=0}^{j}u(i)u(j-i)=\sum_{a,b\ge0}u^\minus_{a+b}(j)=\sum_n(n+1)u^\minus_n(j)=v(j)+u(j)$,
and the term $i=j$ equals $u(j)$. By~\eqref{eq:conv-v} the two forms
of~\eqref{eq:Cj} agree.

We need two facts about sums of plus runs. Write $f=f_\plus$,
$T=T_\plus$, $\alpha=\alpha_\plus$, $\mu=\mu_\plus$ and $S_r=S^\plus_r$ in
Lemmas~\ref{lem:local-bigjump} and~\ref{lem:sums}. Under (LR) the sequence $f$
is regularly varying with index $-\alpha-1$, and Lemma~\ref{lem:rv}(i) gives, for
every $c>0$ and large $n$,
\begin{equation}\label{eq:f-uct}
 \sup_{m\ge cn}\frac{f(m)}{f(n)}\le C_c,\qquad
 \sup_{0\le w\le w_n}\Bigl|\frac{f(n-w)}{f(n)}-1\Bigr|\to0\quad\text{whenever }w_n=o(n).
\end{equation}
By Lemma~\ref{lem:rv}(iv), $\E[X_\plus^2;X_\plus\le y]\le\sum_{n\le y}(2n-1)T(n)\le Cy^2T(y)$
and $\E[X_\plus;X_\plus>y]\to0$. Finally $nT(n)\to0$, since $nT(n)\in\RV_{1-\alpha}$.

\begin{lemma}[uniform local big jump]\label{lem:local-bigjump}
Assume (RV) and (LR) for the plus sign, and fix $\eta\in(0,1)$. Then
\[
 \sup_{1\le r\le n^{1-\eta}}\Bigl|\frac{f^{*r}(n)}{r\,f(n)}-1\Bigr|\to0\qquad(n\to\infty).
\]
\end{lemma}

\begin{proof}
Put $\delta=\eta/2$. By Lemma~\ref{lem:rv}(ii), $rT(n)\le n^{1-\eta}T(n)\le n^{-\delta}$
for large $n$. Fix $\varepsilon\in(0,1/4)$ with $\varepsilon<\delta/(16(1+\delta))$ and
put $y=\varepsilon n$. Let $1\le r\le n^{1-\eta}$ and split the event $\{S_r=n\}$
according to the number of the $r$ runs that exceed $y$.

Exactly one run exceeds $y$. By symmetry this part equals
\[
 E_1=r\sum_{w<n-y}\Prob\bigl(S_{r-1}=w,\ \max_{i<r}X^\plus_i\le y\bigr)\,f(n-w).
\]
Put $w_n=n^{1-\eta/2}$. For $w\le w_n$, \eqref{eq:f-uct} gives $f(n-w)=f(n)(1+o(1))$
uniformly. For $w_n<w<n-y$ we have $n-w>y$ and $f(n-w)\le C_\varepsilon f(n)$
by~\eqref{eq:f-uct}, while $\Prob(S_{r-1}>w_n)\le r\mu/w_n\le\mu n^{-\eta/2}$ by
Markov's inequality. Hence $E_1\le rf(n)(1+o(1))$. Keeping only $w\le w_n$,
\[
 E_1\ge rf(n)(1-o(1))\bigl(1-\Prob(S_{r-1}>w_n)-rT(y)\bigr)=rf(n)(1-o(1)),
\]
since $rT(y)\le C_\varepsilon n^{1-\eta}T(n)\to0$ by Lemma~\ref{lem:rv}(ii).

At least two runs exceed $y$. For two given runs,
$\Prob(S_r=n,\,X^\plus_1>y,\,X^\plus_2>y)\le T(y)\sup_{x>y}f(x)$, by
conditioning on all runs but the second. So this part is at most
$r^2T(y)\sup_{x>y}f(x)\le C_\varepsilon r^2T(n)f(n)$, which is $o(rf(n))$
uniformly, since $rT(n)\to0$.

No run exceeds $y$. For $r=1$ this part vanishes. For $r\ge2$ write
$S_r=S'+S''$, the sums of the first $\lceil r/2\rceil$ and of the remaining
runs. On $\{S_r=n\}$ one of the two halves is at least $n/2$. Bounding the
probability that the other half takes a given value by one,
\[
 E_3=\Prob\bigl(S_r=n,\ \max_iX^\plus_i\le y\bigr)\le2\max_{1\le q\le r}\Prob\bigl(S_q\ge n/2,\ \max_{i\le q}X^\plus_i\le y\bigr).
\]
We apply Lemma~\ref{lem:truncated} with $m=q$, $w=\lfloor\varepsilon n\rfloor+1$ and
$z=n/2-q\mu$. For large $n$ we have $q\mu\le n^{1-\eta}\mu\le n/4$, so $z\ge n/4$, and
$w\le2\varepsilon n$, so $z/(2w)\ge1/(16\varepsilon)$. Also $qT(w)\le rT(\varepsilon n)\le
C_\varepsilon rT(n)<1$ for large $n$, by Lemma~\ref{lem:rv}(ii). Since $a^e$ decreases
in $e$ for $0<a<1$, the lemma gives
\[
 E_3\le2e^{C_{\mathrm K}}\bigl(C_\varepsilon rT(n)\bigr)^{1/(16\varepsilon)} .
\]
By (LR), $f(n)\ge cT(n)/n$. Hence
\[
 \frac{E_3}{rf(n)}\le C'_\varepsilon\,n\,(rT(n))^{1/(16\varepsilon)-1}\le C'_\varepsilon\,n^{1-\delta(1/(16\varepsilon)-1)},
\]
which tends to zero because $\varepsilon<\delta/(16(1+\delta))$ gives
$\delta(1/(16\varepsilon)-1)>1$. The three parts together prove the lemma.
\end{proof}

\begin{lemma}[sums of plus runs]\label{lem:sums}
Assume (RV) for the plus sign. There is a constant $C$ such that, for all $r\ge1$
and real $x>0$,
\[
 \E[S_r;S_r>x]\le r\Bigl(\E[X_\plus;X_\plus>x/2]+\frac{2r\mu^2}{x}\Bigr),
\]
and, if in addition $x\ge4r\mu$,
\[
 \Prob(S_r\ge x)\le C\,rT(x).
\]
\end{lemma}

\begin{proof}
If $S_r\ge x$, either some run is at least $x/2$, or the sum of the runs
truncated at $x/2$ is at least $x$. The truncated sum has mean at most
$r\mu\le x/4$ and variance at most $r\,\E[X_\plus^2;X_\plus<x/2]\le Crx^2T(x)$. By
Chebyshev's inequality the second event has probability at most
$(4/(3x))^2Crx^2T(x)$, and the first at most $rT(x/2)\le CrT(x)$ by
Lemma~\ref{lem:rv}(ii), since $x\ge4r\mu\ge4$. For the other bound, which uses no
condition on $x$, on $\{S_r>x\}$ either $X^\plus_i>x/2$ or $S_r-X^\plus_i>x/2$.
Hence $\E[X^\plus_i;S_r>x]\le\E[X_\plus;X_\plus>x/2]+\mu\Prob(S_r-X^\plus_i>x/2)$, and the
last probability is at most $2r\mu/x$ by Markov's inequality. Sum over $i$.
\end{proof}

\begin{proof}[Proof of Theorem~\ref{thm:second-order}]
We write $u=u_\minus$, $v=v_\minus$ and keep the notation of the two lemmas. By (LR),
$Nf(N)/T(N)\to\alpha$. We use $v(j)\le j\,u(j)\le j$, which holds because
$u^\minus_r(j)=0$ for $r>j$.

(b) Fix $\eta\in(0,1)$, put $J_N=N^{1-\eta}$ and $x_N=N^{1-\eta/2}$, and let $N$ be so
large that $2J_N\le x_N$, $4J_N\mu\le N/2$ and $J_N\le(N/2)^{1-\eta/2}$. Every $o(\cdot)$
below is uniform in $1\le r\le j\le J_N$. In~\eqref{eq:final-explicit} the first sum runs
over $1\le r\le j$, since $u^\minus_0(j)=0$ for $j\ge1$, and the second over $1\le m\le j$
and $0\le r\le j-m$.

The censored final run. Split the expectation in the first sum
of~\eqref{eq:final-explicit} at $S_r=N/2$. For integers $0\le x\le3N/4$ put
\[
 \mathcal T_N(x)=\frac{N}{T(N)}\bigl(T(N-x)-T(N)\bigr)=\frac{N}{T(N)}\sum_{n=N-x}^{N-1}f(n).
\]
By~\eqref{eq:f-uct} and $Nf(N)/T(N)\to\alpha$, $\mathcal T_N(x)=\alpha x(1+o(1))$ uniformly in
$x\le x_N$, and $\mathcal T_N(x)\le Cx$ for $x\le3N/4$, since $f(n)\le C_{1/4}f(N)$ for
$n\ge N/4$. Since $j\le N/4$,
\[
 \E\bigl[T(N-j-S_r);S_r\le N/2\bigr]=T(N)\Bigl(\Prob(S_r\le N/2)+\frac1N\E\bigl[\mathcal T_N(j+S_r);S_r\le N/2\bigr]\Bigr).
\]
On $\{S_r\le x_N/2\}$ we have $j+S_r\le x_N$. With $\E[j+S_r]=j+r\mu$ this gives
\[
 \bigl|\E[\mathcal T_N(j+S_r);S_r\le N/2]-\alpha(j+r\mu)\bigr|
 \le o(j+r)+C\,\E[j+S_r;S_r>x_N/2].
\]
By Markov's inequality and Lemma~\ref{lem:sums},
$\E[j+S_r;S_r>x_N/2]\le2jr\mu/x_N+r\bigl(\E[X_\plus;X_\plus>x_N/4]+4r\mu^2/x_N\bigr)=o(j)$,
since $r\le j\le J_N=o(x_N)$. By Lemma~\ref{lem:sums} and Lemma~\ref{lem:rv}(ii),
$\Prob(S_r>N/2)\le CrT(N)$. On the other part,
\[
 \E\bigl[T(N-j-S_r);N/2<S_r\le N-j-1\bigr]=\sum_{1\le l<N/2-j}f^{*r}(N-j-l)\,T(l).
\]
Here $N-j-l\ge N/2$ and $r\le J_N\le(N-j-l)^{1-\eta/2}$. So Lemma~\ref{lem:local-bigjump}
with $\eta/2$ in place of $\eta$ and~\eqref{eq:f-uct} give $f^{*r}(N-j-l)=rf(N)(1+o(1))$
for $l\le x_N$, since then $j+l=o(N)$, and $f^{*r}(N-j-l)\le CrT(N)/N$ for all $l$ in the
range. Since $\sum_{l\le x_N}T(l)\to\mu$ and $\sum_{l>x_N}T(l)\to0$, the sum equals
$r\alpha\mu T(N)N^{-1}(1+o(1))$. Together,
\[
 \E\bigl[T(N-j-S_r);S_r\le N-j-1\bigr]=T(N)\Bigl(1-O\bigl(rT(N)\bigr)+\frac{\alpha(j+2r\mu)}{N}+o\Bigl(\frac jN\Bigr)\Bigr).
\]
Multiply by $u^\minus_r(j)$ and sum over $r$. Since $\sum_ru^\minus_r(j)=u(j)\le1$ and
$\sum_rru^\minus_r(j)=v(j)$, the first sum of~\eqref{eq:final-explicit} divided by $T(N)$
equals
\[
 u(j)-O\bigl(v(j)T(N)\bigr)+\frac{\alpha}{N}\bigl(j\,u(j)+2\mu v(j)\bigr)+o\Bigl(\frac{j}{N}\Bigr).
\]
Since $v(j)T(N)\le jT(N)=(j/N)\,NT(N)$ and $NT(N)\to0$, the second term is $o(j/N)$.

The final minus run. Since $r+1\le j\le J_N\le(N-J_N)^{1-\eta/2}$,
Lemma~\ref{lem:local-bigjump} with $\eta/2$ in place of $\eta$ and~\eqref{eq:f-uct} give
$u^\plus_{r+1}(N-j)=f^{*(r+1)}(N-j)=(r+1)\alpha T(N)N^{-1}(1+o(1))$. So the second sum
of~\eqref{eq:final-explicit} equals
\begin{align*}
 &\frac{\alpha T(N)}{N}(1+o(1))\sum_{m=1}^{j}T_\minus(m)\sum_r(r+1)u^\minus_r(j-m)\\
 &\qquad=\frac{\alpha T(N)}{N}(1+o(1))\sum_{m=1}^{j}T_\minus(m)\bigl(u(j-m)+v(j-m)\bigr),
\end{align*}
with $v(0)=0$. By~\eqref{eq:cover} and~\eqref{eq:cover-v} the last sum equals
$1+\sum_{i=1}^{j-1}u(i)=\sum_{i=0}^{j-1}u(i)$, which is at most $j$. Adding the two
contributions gives $u(j)+C(j)/N+o(j/N)$ with the first form of~\eqref{eq:Cj}, uniformly
in $1\le j\le J_N$. This is (b).

(a) For $j=0$, \eqref{eq:final-explicit} gives $P_N(N)=T(N)$ exactly, and $C(0)=0$. For
fixed $j\ge1$, part (b) gives an error $o(j/N)=o(1/N)$.

The limit. By~\eqref{eq:renewal-thm} for the minus sign, $u(j)\to1/\mu_\minus$.
By~\eqref{eq:conv-v} and Ces\`aro's theorem, $v(j)/j=j^{-1}\sum_{i<j}u(i)u(j-i)\to\mu_\minus^{-2}$,
and $j^{-1}\sum_{i<j}u(i)\to1/\mu_\minus$. Hence
$C(j)/j\to\alpha_\plus(2/\mu_\minus+2\mu_\plus/\mu_\minus^2)=2\alpha_\plus s/\mu_\minus^2$.
\end{proof}

\subsection{The antimode}\label{sec:antimode}

Two classical facts about the renewal mass function enter. The first is the
second-order renewal theorem. Under (RV) and (SPAN) for the minus sign,
\begin{equation}\label{eq:second-renewal}
 u_\minus(j)-\frac1{\mu_\minus}\sim\frac{j\,T_\minus(j)}{\mu_\minus^2(\alpha_\minus-1)}\qquad(j\to\infty).
\end{equation}
This is Lemma~4 of Frenk~\cite{frenk1982}, stated there for an aperiodic
recurrent event whose tail $\Prob(X>n)\sim T(n)$ is regularly varying with index
$-\alpha$, $\alpha>1$ (see also~\cite{rogozin1973} and~\cite[Theorem~1]{doney2023}).
Section~\ref{sec:verify} reports the check. The second fact gives a simple
condition for the hypothesis of the next theorem.

\begin{lemma}[log-convex run laws]\label{lem:logconvex}
Assume (RV) and (SPAN) for the minus sign. If $f_\minus(n)>0$ for all $n\ge1$
and $f_\minus(n-1)f_\minus(n+1)>f_\minus(n)^2$ for all $n\ge2$, then $u_\minus$ is
strictly decreasing and $u_\minus(j)>1/\mu_\minus$ for every $j\ge0$.
\end{lemma}

\begin{proof}
By Theorem~1 of de Bruijn and Erd\H{o}s~\cite{debruijnerdos1953}, the
renewal sequence of a strictly log-convex law is strictly log-convex, that
is, $u_\minus(m-1)u_\minus(m+1)>u_\minus(m)^2$ for $m\ge1$. Hence the ratios
$u_\minus(m+1)/u_\minus(m)$ increase strictly in $m$. If one of them were at
least $1$, all later ones would exceed some $q>1$, and $u_\minus$ would grow
geometrically, which is impossible since $u_\minus\le1$. So every ratio is
smaller than $1$, $u_\minus$ is strictly decreasing, and it decreases to its
limit $1/\mu_\minus$ given by~\eqref{eq:renewal-thm}.
\end{proof}

\begin{introtheorem}[the antimode]\label{thm:antimode}
Assume (RV) and (SPAN) for both signs and (LR) for the plus sign, and let the
start be fresh with sign $\plus$. Let $j^*_N$ be a minimizer of
$j\mapsto P_N(N-j)$ over $0\le j\le(1-p)N$.
\begin{enumerate}[(a)]
\item If $u_\minus(j)>1/\mu_\minus$ for every $j\ge1$, then $j^*_N\to\infty$,
$j^*_N=o(N)$, and
\[
 T_\minus(j^*_N)\sim\frac{2\alpha_\plus s}{N} .
\]
Moreover $j^*_N/i_N\to1$ for every choice of minimizers $i_N$ of
$j\mapsto u_\minus(j)+C(j)/N$ over $j\ge0$. So $j^*_N$ is of the order of
$b_{\minus,N}$, the stable scale of the minus runs, and
$j^*_N\asymp N^{1/\alpha_\minus}$ up to a slowly varying factor.
\item If $\inf_{j\ge1}u_\minus(j)<1/\mu_\minus$, then $j^*_N$ is bounded and
every value that $j^*_N$ takes infinitely often is a minimizer of $u_\minus$.
\end{enumerate}
\end{introtheorem}

The borderline case, in which $u_\minus\ge1/\mu_\minus$ with equality at some
$j\ge1$, is not covered. By Lemma~\ref{lem:logconvex} part (a) applies to every
strictly log-convex law of the minus runs, for instance the zeta law
$f(n)=n^{-\alpha-1}/\zeta(\alpha+1)$.

\begin{proof}
Write $R_N(j)=P_N(N-j)/T_\plus(N)$, $u=u_\minus$ and
$c_{\mathrm a}=2\alpha_\plus s/\mu_\minus^2$. In this proof $\mu_\minus$ and $\mu_\plus$ are
always written with their signs. We first prove two lower bounds on $R_N$ and then
treat (a) and (b).

Lower bound in the layer. Let $1\le j\le N/(8\mu_\plus)$. In the first sum
of~\eqref{eq:final-explicit}, $T_\plus(N-j-S^\plus_r)\ge T_\plus(N-j)$, and
$\Prob(S^\plus_r\le N-j-1)\ge1-\Prob(S^\plus_r\ge N/2)\ge1-CrT_\plus(N)$ by
Lemma~\ref{lem:sums}, since $4r\mu_\plus\le N/2$. Dropping the second sum,
\begin{equation}\label{eq:anti-lower}
 R_N(j)\ge\frac{T_\plus(N-j)}{T_\plus(N)}\bigl(u(j)-C\,v(j)\,T_\plus(N)\bigr)
 \ge\frac{T_\plus(N-j)}{T_\plus(N)}\bigl(u(j)-C\,j\,T_\plus(N)\bigr).
\end{equation}
By (LR), $f_\plus(n)\ge\frac34\alpha_\plus T_\plus(n)/n\ge\frac34\alpha_\plus T_\plus(N)/N$
for $N/2\le n\le N$ and large $N$, because $T_\plus(n)/n$ is nonincreasing. Hence
$T_\plus(N-j)/T_\plus(N)\ge1+\frac34\alpha_\plus j/N$ for $j\le N/2$. We use
\eqref{eq:anti-lower} below only for $j\le\delta N$ with $\delta\le1/(8\mu_\plus)$ and
on ranges where $u(j)$ is bounded below by a positive constant. There
$CjT_\plus(N)\le C\delta NT_\plus(N)\to0$, so the second factor is positive for large
$N$. Since the first factor is at least $1$, we may then drop it or expand the
product.

Lower bound far from the layer. Fix $\delta\in(0,1-p)$. We claim that
\begin{equation}\label{eq:anti-far}
 \liminf_N\ \min_{\delta N\le j\le(1-p)N}R_N(j)\ge\frac{(1-\delta)^{-\alpha_\plus}}{\mu_\minus}.
\end{equation}
Let $\varepsilon>0$. By Corollary~\ref{cor:bigjump-lower} there is $K$ such that
$P_N(k)\ge(1-\varepsilon)G^{(\plus)}_N(k)$ for $pN+KB_N\le k\le N$ and large $N$.
For $k=N-j$ with $j\ge\delta N$ we have $L_k=N-j/(1-p)\le(1-\delta)N$, so
$G^{(\plus)}_N(k)\ge T_\plus(\lfloor L_k\rfloor)u(j)\ge T_\plus(\lfloor(1-\delta)N\rfloor)u(j)$,
and $T_\plus(\lfloor(1-\delta)N\rfloor)/T_\plus(N)\to(1-\delta)^{-\alpha_\plus}$ by (RV).
For $pN\le k<pN+KB_N$, Theorem~\ref{thm:bulk}(c) and Lemma~\ref{lem:profile}(d)
give $B_NP_N(k)\ge\frac12\inf_N\inf_{[-K,K]}h_N>0$ for large $N$, so
$R_N(N-k)\ge c_K/(B_NT_\plus(N))\to\infty$, since $B_NT_\plus(N)\to0$ by
Lemma~\ref{lem:rv}(v). As $\varepsilon$ is arbitrary and $u(j)\to1/\mu_\minus$, the claim
follows.

(b) Let $m_*=\min_ju(j)<1/\mu_\minus$, which is attained on a finite set because
$u(j)\to1/\mu_\minus$, and let $\varepsilon_*=(1/\mu_\minus-m_*)/3$. Choose $J$ with
$u(j)\ge1/\mu_\minus-\varepsilon_*>0$ for $j\ge J$, and let $\delta=\min\{1/(8\mu_\plus),(1-p)/2\}$. For
$J\le j\le\delta N$, \eqref{eq:anti-lower} gives
$R_N(j)\ge u(j)-CjT_\plus(N)\ge1/\mu_\minus-\varepsilon_*-C\delta NT_\plus(N)$, which exceeds
$m_*+\varepsilon_*$ for large $N$, since $NT_\plus(N)\to0$. For $j\ge\delta N$,
\eqref{eq:anti-far} gives $R_N(j)\ge m_*+\varepsilon_*$ for large $N$. For $j<J$,
Theorem~\ref{thm:second-order}(a) gives $R_N(j)\to u(j)$. So for large $N$ every
minimizer is smaller than $J$, and a value $j'$ taken infinitely often with
$u(j')>m_*$ would give $R_N(j')>R_N(j_0)$ for large $N$, where $u(j_0)=m_*$.

(a) Let $\Lambda_N=a_\minus(N)$, so that $NT_\minus(\Lambda_N)\to1$ and
$\Lambda_N\in\RV_{1/\alpha_\minus}$. Put $w(j)=u(j)-1/\mu_\minus$ and
$c_{\mathrm r}=1/(\mu_\minus^2(\alpha_\minus-1))$, so $w(j)\sim c_{\mathrm r}jT_\minus(j)$
by~\eqref{eq:second-renewal}, and $w(j)>0$ for $j\ge1$. Let
\[
 \Upsilon(x)=c_{\mathrm r}x^{1-\alpha_\minus}+c_{\mathrm a}x,\qquad x>0 .
\]
The function $\Upsilon$ is strictly convex with a unique minimizer $x_*$, given by
$c_{\mathrm r}(\alpha_\minus-1)x_*^{-\alpha_\minus}=c_{\mathrm a}$, that is,
$x_*^{-\alpha_\minus}=c_{\mathrm a}\mu_\minus^2=2\alpha_\plus s$.

We fix the constants first. They depend only on $c_{\mathrm r}$, $c_{\mathrm a}$, $\alpha_\pm$ and
$\mu_\pm$. Choose $\varepsilon\in(0,x_*)$ so small that
$\frac{c_{\mathrm r}}5\varepsilon^{-(\alpha_\minus-1)/2}>\Upsilon(x_*)+2$, and choose $K>x_*$ so large
that $\alpha_\plus K/(2\mu_\minus)>\Upsilon(x_*)+1$. So $\varepsilon<x_*<K$.

Step 1 (the window). By~\eqref{eq:second-renewal} and
Lemma~\ref{lem:rv}(i),
$(N/\Lambda_N)\,w(j)=c_{\mathrm r}(j/\Lambda_N)^{1-\alpha_\minus}(1+o(1))$ uniformly in
$\varepsilon\Lambda_N\le j\le K\Lambda_N$, because
$NT_\minus(j)=NT_\minus(\Lambda_N)\cdot T_\minus(j)/T_\minus(\Lambda_N)$. By
Theorem~\ref{thm:second-order}, $C(j)=c_{\mathrm a}j(1+o(1))$ as $j\to\infty$. Choose
$\eta>0$ with $1-\eta>1/\alpha_\minus$, so that $K\Lambda_N\le N^{1-\eta}$ for large
$N$. Theorem~\ref{thm:second-order}(b) then gives
\begin{equation}\label{eq:anti-window}
 \frac{N}{\Lambda_N}\Bigl(R_N(j)-\frac1{\mu_\minus}\Bigr)=\Upsilon\Bigl(\frac{j}{\Lambda_N}\Bigr)+o(1)
 \quad\text{and}\quad
 \frac{N}{\Lambda_N}\Bigl(u(j)+\frac{C(j)}N-\frac1{\mu_\minus}\Bigr)=\Upsilon\Bigl(\frac{j}{\Lambda_N}\Bigr)+o(1),
\end{equation}
uniformly in $\varepsilon\Lambda_N\le j\le K\Lambda_N$. Since $\varepsilon<x_*<K$, in
particular the minimum of $R_N$ is at most $1/\mu_\minus+(\Lambda_N/N)(\Upsilon(x_*)+o(1))$.

Step 2 (small $j$). We have $R_N(0)=1$, and $u(0)+C(0)/N=1$ for the comparison
function. Since $\mu_\minus>1$ and $\Lambda_N/N\to0$, both exceed
$1/\mu_\minus+(\Lambda_N/N)(\Upsilon(x_*)+1)$ for large $N$. Let
$1\le j\le\varepsilon\Lambda_N$. Then $u(j)>1/\mu_\minus$, and~\eqref{eq:anti-lower} gives
$R_N(j)\ge u(j)-C\varepsilon\Lambda_NT_\plus(N)=1/\mu_\minus+w(j)-o(\Lambda_N/N)$. Choose
$j_0$ with $w(j)\ge\frac{c_{\mathrm r}}2jT_\minus(j)$ for $j\ge j_0$. By
Lemma~\ref{lem:rv}(ii) applied to $jT_\minus(j)\in\RV_{1-\alpha_\minus}$,
for $j_0\le j\le\varepsilon\Lambda_N$ and $j_0$ large,
\[
 \frac{jT_\minus(j)}{\Lambda_NT_\minus(\Lambda_N)}\ge\frac12\Bigl(\frac{j}{\Lambda_N}\Bigr)^{-(\alpha_\minus-1)/2}\ge\frac12\varepsilon^{-(\alpha_\minus-1)/2}.
\]
Since $\Lambda_NT_\minus(\Lambda_N)\sim\Lambda_N/N$, we get
$w(j)\ge\frac{c_{\mathrm r}}5\varepsilon^{-(\alpha_\minus-1)/2}\Lambda_N/N$ for large $N$.
For $1\le j<j_0$, $w(j)\ge\min_{i<j_0}w(i)>0$, which is much larger than
$\Lambda_N/N$. By the choice of $\varepsilon$,
$R_N(j)>1/\mu_\minus+(\Lambda_N/N)(\Upsilon(x_*)+1)$ for $1\le j\le\varepsilon\Lambda_N$ and
large $N$. The same bound holds for $u(j)+C(j)/N\ge1/\mu_\minus+w(j)$.

Step 3 (intermediate $j$). Let $\delta=\min\{1/(8\mu_\plus),(1-p)/2\}$ and
$K\Lambda_N\le j\le\delta N$. By~\eqref{eq:anti-lower} and $u(j)>1/\mu_\minus$,
\[
 R_N(j)\ge\Bigl(1+\frac{3\alpha_\plus j}{4N}\Bigr)\Bigl(\frac1{\mu_\minus}-CjT_\plus(N)\Bigr)
 \ge\frac1{\mu_\minus}+\frac{\alpha_\plus j}{2\mu_\minus N}\ge\frac1{\mu_\minus}+\frac{\alpha_\plus K}{2\mu_\minus}\cdot\frac{\Lambda_N}{N}
\]
for large $N$, since $NT_\plus(N)\to0$. By the choice of $K$ this exceeds
$1/\mu_\minus+(\Lambda_N/N)(\Upsilon(x_*)+1)$. For the comparison function,
$C(j)\ge\alpha_\plus j\,u(j)\ge\alpha_\plus j/\mu_\minus$, so
$u(j)+C(j)/N\ge1/\mu_\minus+\alpha_\plus j/(\mu_\minus N)$ for every $j$, and the same bound
holds for all $j\ge K\Lambda_N$.

Step 4 (large $j$). By~\eqref{eq:anti-far}, $R_N(j)\ge(1+\varepsilon_\delta)/\mu_\minus$ for
$\delta N\le j\le(1-p)N$ and large $N$, with $\varepsilon_\delta=((1-\delta)^{-\alpha_\plus}-1)/2>0$, while
$\Lambda_N/N\to0$.

Conclusion. By Steps 2 to 4 every minimizer of $R_N$, and every minimizer of
$u(j)+C(j)/N$, lies in $[\varepsilon\Lambda_N,K\Lambda_N]$ for large $N$. Fix
$\eta'>0$. Since $\Upsilon$ is continuous with a unique minimizer, $\Upsilon\ge\Upsilon(x_*)+\varepsilon'$
on $[\varepsilon,K]\setminus(x_*-\eta',x_*+\eta')$ for some $\varepsilon'>0$, and
\eqref{eq:anti-window} shows that both minimizers lie in
$((x_*-\eta')\Lambda_N,(x_*+\eta')\Lambda_N)$ for large $N$. Hence
$j^*_N/\Lambda_N\to x_*$ and $i_N/\Lambda_N\to x_*$, so $j^*_N/i_N\to1$. By (RV),
$NT_\minus(j^*_N)=NT_\minus(\Lambda_N)\cdot T_\minus(j^*_N)/T_\minus(\Lambda_N)\to x_*^{-\alpha_\minus}=2\alpha_\plus s$.
Finally $\Lambda_N=a_\minus(N)$ and $b_{\minus,N}=a_\minus(N/s)$ differ by a
bounded factor, and $a_\minus\in\RV_{1/\alpha_\minus}$.
\end{proof}

The leading law $T_\minus(j^*_N)\sim2\alpha_\plus s/N$ is approached slowly,
because $C(j)/j$ converges slowly, and the minimizer of $u_\minus(j)+C(j)/N$ is the
better prediction at moderate $N$. Section~\ref{sec:verify} reports the check.
